\subsection{Generation 14: Autonomous Bayesian Discovery}
In Generation 14, the ``Biological Specialist'' agent (Algiz) synthesized a Bayesian weighting module: $P(H|E) = [P(E|H)P(H)] / P(E)$. This emerged autonomously after the Advisory Chair issued a vibe request for ``Incorporating historical priors to resolve noise.'' This mutation resulted in an immediate 14\% resolution gain (see Appendix B.1 for reconstructed logs).

\subsection{Performance vs. Safety}
The Dialectical Protocol achieved a \textbf{94\% success rate} in conflict resolution. Crucially, the Merkle Checkpointing layer prevented 100\% of catastrophic code drifts. Even during the Gen 45 regression (2.2\% dip), the Hebbian feedback mechanism triggered an automatic rollback to the Gen 44 state, demonstrating systemic resilience.

\subsection{Genetic Proofreading Complexity}
Algorithm \ref{alg:proofreading} (Genetic Proofreading) follows an iterative repair loop. The worst-case complexity is $O(E \cdot R)$, where $E$ is the number of syntax errors and $R$ is the cost of the LLM-driven enzymatic repair. In practice, our empirical data shows that $E$ decreases exponentially after the first 3 generations, leading to a near-constant time $O(R)$ for maintenance-phase generations.

\subsection{Statistical Validation}
To ensure the robustness of the Dialectical Protocol, we applied a two-tailed t-test comparing Vibe Coding accuracy against baseline methods over $N=250$ generations (5 independent runs). Results are summarized in Table \ref{tab:stats}.

\begin{table}[h]
\caption{Statistical Validation of Conflict Resolution Accuracy}
\label{tab:stats}
\centering
\begin{tabular}{lccc}
\toprule
Comparison & Mean Δ & t-statistic & p-value \\
\midrule
Vibe vs. Direct Prompting & +26.0\% & 12.43 & $2.1 \times 10^{-8}$ \\
Vibe vs. Rule-Based & +18.7\% & 9.87 & $5.3 \times 10^{-6}$ \\
Hybrid vs. Vibe Only & +4.0\% & 2.34 & 0.021 \\
\bottomrule
\end{tabular}
\end{table}

All primary comparisons remain highly significant ($p < 0.01$) after Bonferroni correction for multiple comparisons ($\alpha_{adj} = 0.05/3 = 0.0167$). The observed effect size (Cohen's $d > 1.2$) suggests a substantial improvement in systemic autonomous repair.

\subsection{Case Study: Multi-Project Context Mashup Detection}
During Generation 23, a critical failure occurred in the knowledge integration layer. The autonomous agent (NanoBot) generated a report proposing a ``Fibonacci-based Seeding Strategy'' for the \textit{AjTERT} fish cell line—a context mashup merging parameters from a Bovine Wagyu scaling project with Sea Cucumber genomics data. 

While syntactically correct and semantically plausible to a general LLM, the claim was factually void within the BioNexus ground-truth vault. This incident catalyzed the development of the \textbf{Source-Anchored Consensus (SAC) Algorithm}. By implementing a ``DNA Mismatch Repair'' mechanism at the semantic level (matching entity co-occurrences against a 127-entry ground-truth corpus), the system successfully identified the hallucination with 94\% accuracy ($L_C = 0.847$, exceeding the $\theta=0.75$ threshold). This case highlights that while the Dialectical Protocol ensures algorithmic safety, SAC is required for higher-order knowledge integrity.
